% ============ 03. Ciencia de la ciencia y evaluación responsable ============
\section{Ciencia de la Ciencia (SciSci) y evaluación responsable}

\subsection{Qué dice la literatura}

La revisión fundacional de \citet{fortunato2018science} consolidó la
\textbf{Ciencia de la Ciencia (SciSci)}: el estudio empírico del sistema
científico con datos masivos (publicaciones, registros administrativos, redes)
y métodos de la física estadística, las redes complejas y el aprendizaje
automático. Sus líneas maduras — dinámicas de carrera, productividad,
colaboración, movilidad, desigualdad y estructura institucional — son
exactamente los ejes que aborda este repositorio. La síntesis de
\citet{zeng2021science} y el libro de \citet{wang2021sciencebook} ordenan esos
resultados; \citet{clauset2017data} mostró que los datos bibliométricos y de
registro permiten predicciones cuantitativas sobre carreras individuales
(edad de pico, probabilidad de alta productividad), con implicaciones directas
para un observatorio de capital humano científico.

En paralelo, el movimiento de \textbf{evaluación responsable} fijó los límites
de uso de los indicadores: la Declaración de San Francisco
\citep{dora2012} (2012), los \textbf{Principios de Leiden}
\citep{hicks2015leiden} (2015) y el informe \emph{Metric Tide}
\citep{wilsdon2015metric} (2015) coinciden en que (i) los indicadores
cuantitativos apoyan, no sustituyen, el juicio; (ii) deben ser robustos al
contexto de campo; y (iii) sus definiciones y límites deben ser transparentes.
El \textbf{Acuerdo de CoARA} \citep{coara2022} (2022) lleva esta agenda a la
práctica institucional en Europa (reforma de la evaluación de la
investigación), y existe debate activo sobre su aplicación en América Latina,
donde se ha señalado que ciertos indicadores "desestimulan" formas legítimas de
producción de conocimiento (por ejemplo, los libros en ciencias sociales y
humanidades; véase la crítica de CLACSO citada en la sección 13.2).

\subsection{Relación con este repositorio}

El proyecto ya adopta el espíritu SciSci (patrón \emph{Pregunta →
Hallazgo → Caveat} y hallazgos sobre retención, concentración, brechas y
productividad) y su crítica del sistema de reconocimiento colombiano está
alineada con la agenda de métricas responsables. La auditoría verificó, sin
embargo, dos brechas frente a esta literatura: (a) \textbf{no existe un marco
formal de indicadores} (definición operativa, fórmula, fuente y límites de cada
métrica), condición mínima del enfoque de Leiden/CoARA; y (b) la inferencia es
descriptiva, sin modelos predictivos ni de robustez, cuando SciSci ofrece
herramientas maduras para ello.

\subsection{Mejoras recomendadas}

\begin{enumerate}
  \item Publicar un \textbf{marco de indicadores responsable}: para cada métrica
        del observatorio (retención, transición, HHI, brechas, representación),
        declarar fórmula, fuente, población, supuestos y usos permitidos
        (Leiden; DORA; CoARA).
  \item Adoptar \textbf{análisis de robustez} (bootstrap, sensibilidad a
        exclusiones) como estándar de publicación de cada hallazgo
        (recomendado por \citet{hicks2015leiden} para evitar sobreinterpretación).
  \item Explorar \textbf{modelos predictivos de carreras} (estilo
        \citet{clauset2017data} y \citet{wang2021sciencebook}) como siguiente
        fase del observatorio: probabilidad de promoción/deserción por cohorte,
        edad y área — insumo de política de alto valor.
\end{enumerate}
